\documentclass[11pt]{article}
\usepackage[a4paper,margin=18mm]{geometry}
\usepackage{fontspec}
\setmainfont{Latin Modern Roman}
\usepackage{amsmath,amssymb,amsthm,amscd,graphicx,color}
\usepackage{xurl}
\usepackage[unicode,hidelinks]{hyperref}
\allowdisplaybreaks
\setlength\emergencystretch{3em}

\makeatletter
\newtheorem{theorem}{Theorem}
\newtheorem*{theorem*}{Theorem}
\newtheorem{lemma}{Lemma}
\newtheorem*{lemma*}{Lemma}
\newtheorem{corollary}{Corollary}
\newtheorem*{corollary*}{Corollary}
\newtheorem{proposition}{Proposition}
\newtheorem*{proposition*}{Proposition}
\newtheorem{conjecture}{Conjecture}
\newtheorem*{conjecture*}{Conjecture}
{\theoremstyle{definition} \newtheorem{definition}{Definition}
\newtheorem*{definition*}{Definition}
\newtheorem{example}{Example}
\newtheorem*{example*}{Example}
\newtheorem{remark}{Remark}
\newtheorem*{remark*}{Remark}
\newtheorem{note}{Note}
\newtheorem*{note*}{Note}
}
\makeatother

\DeclareMathAlphabet{\SigmaLegacyBold}{OT1}{cmr}{bx}{n}\renewcommand{\bf}{\ifmmode\expandafter\SigmaLegacyBold\else\expandafter\bfseries\fi}
\makeatletter\newlabel{secth1}{{2}{}{Original source locator}{}{}}\makeatother
\begin{document}
\begin{center}\Large Small-data wave operators for Moyal nonlinear Schrodinger evolution in dimension two d with alpha greater than two d and real polynomial interaction without constant term, and without linear term when d is one or two\end{center}
\noindent\textbf{Stable ID:} 2782\par
\noindent\textbf{Authors:} Bergfinnur Durhuus; Victor Gayral\par
\noindent\textbf{Original work:} The Scattering Problem for a Noncommutative Nonlinear Schrödinger Equation\par
\noindent\textbf{Version:} Official SIGMA 6 (2010) article 046, DOI 10.3842/SIGMA.2010.046; journal PDF and native TeX acquired 2026-10-01, exact bytes pinned by SHA256\par
\noindent\textbf{Selected task scope:} Theorem3 in the source oscillator weighted Hilbert space H\_alpha=\allowbreak{}L\_\{2,alpha\}, b\_mn=\allowbreak{}1+\allowbreak{}EuclideanNorm(m-n), and Sigma\_alpha with norm ||phi||\_\{2,alpha\}+\allowbreak{}sup\_t abs(t)\textasciicircum{}(d/2)||exp(-it Delta\_alpha)phi||\_a. Let d>=\allowbreak{}1, alpha>2d, and F a real polynomial with F(0)=\allowbreak{}0, with zero linear coefficient when d=\allowbreak{}1 or2. There is delta>0 such that scattering data phi\_plus/minus in Sigma\_alpha of norm<delta have unique global continuous solutions of original equations(9)/(10), converging in H\_alpha to free evolution as t goes to the corresponding infinity. If F also has zero linear coefficient, the free-frame difference converges in the full scattering norm. Theta>0 is absorbed into F as stated in the original equation; no large-data general or asymptotic completeness assertion.\par
\noindent\textbf{Difficulty:} H2: Classical dispersive norms are poorly behaved under the Moyal product, so standard small-data NLS scattering cannot be applied without constructing another analytic framework. The author obtains uniform oscillator-matrix propagation bounds by rewriting the heat kernel as Jacobi polynomials and applying a precisely stated uniform weighted bound. New weighted matrix-product and polynomial Lipschitz estimates then verify every hypothesis of the abstract nonlinear scattering theorem, yielding a concrete scattering space and the stated degree restrictions.\par
\noindent\textbf{Editorial boundary note (not author text):} The complete original Theorem3 is retained for the precise oscillator-weighted scattering space, alpha>2d and the stated polynomial degree restrictions. The original Moyal/Weyl/operator formulation, weighted norms, Jacobi heat-kernel rewrite, full uniform propagation estimate, matrix product bounds and all four nonlinear scattering hypotheses are bundled. The cited uniform Jacobi polynomial bound and Reed scattering theorem are named prior prerequisites, not new claims needing recursive proofs. Neither general large-data scattering nor separately unselected diagonal improvements/asymptotic completeness is selected.\par
\noindent\textbf{Source:} \url{https://sigma-journal.com/2010/046/}\quad DOI: \url{https://doi.org/10.3842/SIGMA.2010.046}\par
\noindent\textbf{License and attribution:} Copyright retained by the named authors. Published by SIGMA under CC BY 4.0: \url{https://creativecommons.org/licenses/by/4.0/}. Source: \url{https://sigma-journal.com/about.html}. Adaptation consists of standalone excerpt formatting, cover and numbering restoration; original author language and mathematical text are retained.\par
\noindent\textbf{Typesetting adaptation (not author text):} The legacy math-bold switch is mapped to the standard installed Computer Modern bold math alphabet so the author bold Delta operator retains its glyph and distinction from ordinary Delta under portable XeLaTeX. Author formula tokens are unchanged.\par
\noindent\textbf{External original locator secth1:} Section2 proves separately unselected large-data Cauchy existence. The selected small scattering-data theorem instead closes by the explicitly cited Reed scattering theorem after all new norm/propagation hypotheses are verified.\par
\medskip\hrule\medskip\noindent\textbf{Author text begins below.}\par
\setcounter{section}{0}
\setcounter{subsection}{0}
\setcounter{subsubsection}{0}
\setcounter{equation}{0}
% Original source lines 63--213
\section{Introduction}

The scattering problem for general nonlinear
f\/ield equations has been intensively studied for many years with
considerable progress, a seminal result being the establishment of
asymptotic completeness in energy space for the
nonlinear Schr{\"o}dinger equation in space dimension \mbox{$n\geq 3$} with
interaction $|\varphi|^{p-1}\varphi$, where
$1+4/n\leq p\leq 1+4/(n-2)$, and analogous results for the nonlinear Klein--Gordon
equation~\cite{GiVe}. Still, many important questions remain
unanswered, in particular relating to the existence of wave operators
def\/ined on appropriate subspaces of scattering states and to
asymptotic completeness for more general interactions, although many
partial results have been obtained. We refer to~\cite{Strauss,Ginibre} for an overview.
One source of dif\/f\/iculties encountered can be traced back to the
singular nature of pointwise
multiplication of functions with respect to appropriate Lebesgue or
Sobolev norms. It therefore appears natural to look for f\/ield equations where
this part of the problem can be eliminated while preserving as much
as possible of the remaining structure. One such possibility
is to deform the standard product appropriately, and  our
purpose  in the present article is to exploit this idea.

Deformations of the kind mentioned applied to relativistic f\/ield
equations have turned out to be of relevance for understanding
non-perturbative aspects of string theory, see e.g.~\cite{GMS,HKL,DJN}.
We shall limit ourselves to studying a deformed version of the
nonlinear Schr{\"o}dinger equation in an even number $n=2d$ of space
dimensions. The rather crude Hilbert space
techniques~\cite{Reed,RS} we apply allow us to consider polynomial
interactions only. We hope that an application of more modern techniques
\cite{Caz} will allow a treatment of a larger class of interactions.
Our main purpose will be to demonstrate that, for
the deformed equation,  the Cauchy problem is in a certain sense more
regular as compared to the
classical equation as well as to set up a natural scattering framework,
 including appropriate decay estimates of the free time evolution
and construction of wave operators under rather mild conditions on the
interaction.

The deformed, or
noncommutative, version of the nonlinear Schr{\"o}dinger
equation (NCNLS) in $2d$ space dimensions can be written as
\begin{gather}
\label{NLSch}
\big(i\partial_t-{\Delta}\big)\varphi(x,t) =
\varphi\,\star_\theta F_{\star_\theta}(\varphi^*\star_\theta\varphi)(x,t) ,
\end{gather}
where $\star_\theta$ denotes the Moyal product (see below) of functions of
$2d$ space variables $x$, $F_{\star_\theta}$ denotes a real polynomial
with respect to this product and ${\Delta}=-\sum_{i=1}^{2d}\partial_i^2$ is
the standard Laplacian in~$2d$ variables.
The Moyal product considered here is given by
\begin{gather*}
%\label{mop}
f\star_\theta g(x):=(2\pi)^{-2d}\int e^{-iyz} f\left(x-\frac\Theta2
y\right)  g(x+z)\,d^{2d} y\,d^{2d} z .
\end{gather*}
The constant skew-symmetric $(2d\times2d)$-matrix $\Theta$ is assumed to
be given in the canonical form
\[
\Theta  =  \theta \begin{pmatrix} 0 & I_d\\ -I_d & 0\end{pmatrix} ,
\]
where $I_d$ denotes the $d\times d$ identity matrix and
$\theta>0$ is called the deformation parameter.

Def\/ining
$\varphi_\theta(x,t) = \varphi(\theta^\frac 12 x,\theta t)$,
we have the scaling identity
\[
(\varphi\star_\theta \psi)_{\theta} = \varphi_{\theta}\star_1 \psi_{\theta}.
\]
It follows that $\varphi$ satisf\/ies (\ref{NLSch}) if and only if
\begin{gather*}
%\label{NLSch1}
\big(i\partial_t-\Delta\big)\varphi_\theta(x,t) =
\theta \varphi_\theta \star F_\star(\varphi_\theta^*\star\varphi_\theta)(x,t) ,
\end{gather*}
where we have dropped the subscript on the $\star$-product when $\theta =1$.


The $\star$-product is intimately connected to the so-called Weyl
quantization map $W$. This map associates an operator
$W(f)$ on $L^2(\mathbb R^d)$ to appropriate functions (or distributions) $f$ on $\mathbb R^{2d}$
such that the kernel $K_W(f)$ of $W(f)$ is given by
\begin{gather*}
%\label{W}
K_W(f)(x,y)=(2\pi)^{-d}\int_{\mathbb
R^d}f\left(\frac{x+y}{2},p\right)e^{i(x-y)\cdot p}\,dp  .
\end{gather*}
 It is easily seen that $W$ is an
isomorphism from $L^2(\mathbb R^{2d})$ onto the Hilbert space $\cal H$ of
Hilbert--Schmidt operators on $L^2(\mathbb R^d)$ fulf\/illing
\begin{gather}
\label{Wunitary}
\Vert W(f)\Vert_2 = (2\pi)^{-d/2}\Vert f\Vert_{L^2(\mathbb R^{2d})} ,
\end{gather}
where $\Vert\cdot\Vert_2$ denotes the Hilbert--Schmidt norm. Moreover,
$W$ maps the space ${\cal S}(\mathbb R^{2d})$ onto the space of
operators whose kernel is a Schwartz function and its relation to the
$\star$-product is exhibited by the identity
\[
W(f\star g) = W(f)W(g) .
\]
Further properties of the $\star$-product can be found in
e.g.~\cite{GGBISV}.

By use of $W$, equation~\eqref{NLSch} can be restated (see also \cite{DJN})
as a dif\/ferential equation
\begin{gather}
\label{NLSch2}
i\partial_t\phi - 2 \sum_{k=1}^d[a_k^*,[a_k,\phi]] =\theta \phi F(\phi^*\phi) ,
\end{gather}
 for  the operator-valued function $\phi(t) := W(\varphi_\theta(\cdot,t))$,
where we have introduced the creation and annihilation operators
\[
a_k =\frac{1}{\sqrt 2}(x_k + \partial_k)\qquad\mbox{and}\qquad a_k^* =
\frac{1}{\sqrt 2}(x_k - \partial_k) .
\]
The operator
\[
{\bf \Delta} = 2 \sum_{k=1}^d \mbox{ad}\, a_k^*\;\mbox{ad}\, a_k,
\]
with domain $D({\bf\Delta})$,
is def\/ined in a natural way as a self-adjoint operator on $\cal H$ by the
relations
\begin{gather}
\label{tria}
{\bf \Delta} =  W\Delta W^{-1} ,\qquad D({\bf \Delta})=WD(\Delta) ,
\end{gather}
where $\Delta$ denotes the standard self-adjoint $2d$-dimensional Laplace
operator with maximal domain $D(\Delta)=H_2^2(\mathbb R^{2d})$.

We shall primarily be interested in globally def\/ined \emph{mild}
solutions to the Cauchy problem associated to
equation (\ref{NLSch2}), that is continuous solutions $\phi:\mathbb
R\to{\cal H}$ to the corresponding integral equation
\begin{gather}
\label{NLSch3}
\phi(t) = e^{-i(t-t_0){\bf \Delta}}\phi_0 - i\int_{t_0}^t e^{i(s-t){\bf \Delta}}
\phi(s) F\big(|\phi(s)|^2\big)\,ds .
\end{gather}
This equation is weaker than (\ref{NLSch2}) in the sense that if
$\phi: I\to D({\bf \Delta})$ is a continuously dif\/ferentiable solution to
(\ref{NLSch2}) def\/ined on an interval $I$ containing $t_0$, i.e.\ a
\emph{strong solution} on $I$, then it also fulf\/ills (\ref{NLSch3}).
This latter equation  naturally f\/its into
the standard Hilbert space framework for evolution equations, see
e.g.~\cite{RS,Reed}. The following theorem is proven in Section~\ref{secth1}
below.

\setcounter{section}{1}
\setcounter{subsection}{0}
\setcounter{subsubsection}{0}
\setcounter{equation}{5}
\setcounter{theorem}{1}
% Original source lines 244--387


 In order to formulate our main results on the scattering problem we
need to introduce ap\-propriate Hilbert spaces and auxiliary norms.
By $\vert n\rangle$, $n=(n_1,\dots,n_k)\in\mathbb N_0^d$ we shall denote the
standard orthonormal basis for $L^2(\mathbb R^d)$ consisting of
eigenstates for the $d$-dimensional harmonic oscillator, and by
$(\phi_{mn})$ we denote the matrix representing a bounded operator
$\phi$ with respect to this basis, that is
\[
\phi = \sum_{m,n}\phi_{mn}\vert n\rangle\langle m|,\qquad \phi_{mn}:=\langle n|\phi|m\rangle .
\]
Let
\[
b_{mn} := 1+|m-n| ,\qquad m,n\in \mathbb N_0^d ,
\]
where $|\cdot|$ denotes the Euclidean norm in $\mathbb R^d$. For
exponents $p\geq 1$, $\alpha\in\mathbb R$ and $\phi$ as above, we def\/ine
the norms $\Vert\cdot \Vert_{p,\alpha}$  by
\begin{gather}
\label{alphanorms}
\Vert\phi\Vert_{p,\alpha}^p := \sum_{m,n}b_{mn}^{\alpha p}|\phi_{mn}|^p ,
\quad p<\infty, \qquad\mbox{and}\qquad
\Vert\phi\Vert_{\infty,\alpha}:=\sup_{m,n}\{b_{mn}^\alpha|\phi_{mn}|\} ,
\end{gather}
and we let ${\cal L}_{p,\alpha}$ denote the space of operators
on $L^2(\mathbb R^d)$ for which
$\Vert\cdot \Vert_{p,\alpha}$ is f\/inite. We note that
$\mathcal H_{\alpha}:=\mathcal L_{2,\alpha}$ is a  Hilbert space with
scalar product
\[
\langle\phi,\psi\rangle_\alpha:=\sum_{m,n}b_{mn}^{2\alpha} \overline{\phi_{mn}}
 \psi_{mn} ,
\]
and, in particular, ${\cal H}_{0}$ equals  the space $\cal H$ of
Hilbert--Schmidt operators on $L^2(\mathbb R^d)$ with $\Vert\cdot\Vert_{2,0} =
\Vert\cdot\Vert_2$.
Clearly, the operator $U_\alpha:\mathcal H_\alpha\to \mathcal H$ def\/ined by
\begin{gather}\label{Ualpha}
(U_\alpha\phi)_{mn}=
b^{\alpha}_{mn}\phi_{mn} ,\qquad \phi\in\mathcal H_\alpha ,
\end{gather}
is unitary, and the operator
${\bf\Delta}_\alpha := U_\alpha^*\, {\bf \Delta}\,U_\alpha$, with domain
$D({\bf\Delta}_\alpha):=U_\alpha^* D({\bf \Delta})$,
 is self-adjoint on $\mathcal H_\alpha$.

We use the notation $\Vert\cdot\Vert_{\rm op}$ for
the standard operator norm and def\/ine the norm $\Vert\cdot\Vert_a$ by
\[
\Vert\phi\Vert_a := \Vert\widetilde\phi\Vert_{\rm op} ,
\]
whenever the operator
\[
 \widetilde\phi := \sum_{m,n}|\phi_{mn}|\, |n\rangle\langle m|
\]
 is bounded.


With this notations we have the following decay estimate for the free
propagation:
\begin{theorem}\label{th2} For $\alpha\geq 0$ and  $\beta > d$ there exists a constant
$c_\beta>0$ such that
\begin{gather}\label{freedecay}
\Vert\, e^{-it{\bf\Delta}_\alpha}\phi \Vert_a  \leq  c_{\beta}(1+
|t|)^{-\frac{d}{2}} \Vert\phi\Vert_{1,\beta} ,
\qquad \phi\in {\cal H}_\alpha\cap {\cal L}_{1,\beta} .
\end{gather}
\end{theorem}

This estimate should be compared to the standard one used for the
classical nonlinear wave equations with the $L^\infty$-norm on the
left and a $L^1$-Sobolev norm on the right, applied to
functions of $2d$ space variables. In this case, one obtains rather
trivially a decay exponent~$d$ instead of~$d/2$. However, those norms
are not well behaved with respect to the $\star$-product and cannot be
used for our purposes. Instead, we have to work a bit harder to establish
(\ref{freedecay}) using uniform estimates on the classical Jacobi
polynomials. This is accomplished in Section~\ref{secth2}.



Our last main result concerns the existence of wave operators
def\/ined on a scattering subspace $\Sigma_\alpha\subset \mathcal H_\alpha$. In order to def\/ine
the latter we introduce, for f\/ixed $\alpha\geq 0$, the scattering norm
of $\phi\in{\cal H}_\alpha$ by
\begin{gather*}%\label{scatnorm}
|||\phi|||_{\alpha}:=
\Vert\phi\Vert_{2,\alpha}+\sup_{t\in\mathbb R} |t|^{\frac d2}
\Vert e^{-it{\bf\Delta}_\alpha}\phi\Vert_a
\end{gather*}
and set
\begin{gather*}
%\label{sigmascat}
\Sigma_{\alpha}:=\left\{\phi\in\mathcal H_\alpha\,|\;|||\phi|||_{\alpha}<\infty\right\}\,.
\end{gather*}
From \eqref{alphanorms}, we immediately see that
$\mathcal L_{p,\alpha}\subset\mathcal L_{q,\beta}$, if $p\leq q$ and
$\alpha p\geq\beta q$, so that Theorem \ref{th2} gives, in particular, the
inclusion $\mathcal L_{1,\beta}\subset\Sigma_\alpha$ if $\beta> \max\{d, 2\alpha\}$.

The relevant integral equations to solve in a scattering situation
correspond to initial/f\/inal data $\phi_\pm$ at $t=\pm\infty$ and
take the form
\begin{gather}
\label{NLSch4-}
\phi(t)=e^{-it{{\bf\Delta}_\alpha}}\phi_--i\int_{-\infty}^t
e^{i{(s-t){\bf\Delta}_\alpha}}\phi(s) F\big(|\phi(s)|^2\big)\,ds ,
\end{gather}
and
\begin{gather}
\label{NLSch4+}
\phi(t)=e^{-it{{\bf\Delta}_\alpha}}\phi_++i\int^{+\infty}_t
e^{i{(s-t){\bf\Delta}_\alpha}}\phi(s) F\big(|\phi(s)|^2\big)\,ds ,
\end{gather}
respectively, where $F$ has been redef\/ined to include $\theta$.
 {\it We assume here for simplicity that the polynomial $F$
has no constant term}, since such a term could trivially be incorporated by
subtracting it from ${\bf\Delta}_\alpha$ in the preceding formulas.   The
f\/irst problem
then is to determine spaces of initial/f\/inal data $\phi_\pm$, such that the
equations above have a unique global solution and such that these
solutions behave as free solutions for $t\to \pm\infty$.
To this end we establish the following in Section~\ref{secth3}.

\begin{theorem}\label{th3}
Let $\alpha>2d$ and assume that the polynomial $F$ has no constant
term and, in addition, no linear term
if $d=1$ or $d=2$.

Then there exists $\delta>0$ such that for every
 $\phi_\pm\in\Sigma_{\alpha}$ with $|||\phi_\pm|||_{\alpha}<\delta$, the
equations~\eqref{NLSch4-} and~\eqref{NLSch4+} have unique globally
defined continuous solutions $\phi^\pm:\mathbb R\to \Sigma_{\alpha}$ fulfilling
\begin{gather}\label{solsmall1}
\Vert\phi^\pm(t)- e^{-it{{\bf\Delta}_\alpha}}\phi_\pm\Vert_{2,\alpha}  \to
0,\qquad\mbox{for}\ \ t\to\pm\infty .
\end{gather}
If, furthermore, $F$ has no linear term, then we have
\begin{gather}\label{solsmall2}
 |||e^{it{{\bf\Delta}_\alpha}}\phi^\pm(t)- \phi_\pm |||_{\alpha}  \to
0,\qquad\mbox{for}\  \ t\to\pm\infty .
\end{gather}
\end{theorem}
\setcounter{section}{2}
\setcounter{subsection}{0}
\setcounter{subsubsection}{0}
\setcounter{equation}{17}
% Original source lines 527--860
\section{Some norm inequalities}
\label{secframe}

In preparation for the proofs of Theorems \ref{th2} and \ref{th3} we
collect in this section a few useful lemmas.

\begin{lemma}
\label{Bbounded}
For $\alpha<-d$, the matrix
$\{b_{mn}^\alpha\}$ represents a bounded operator $B_\alpha$ on
$L^2(\mathbb R^d)$ with respect to the basis $\{|n\rangle\}$.
\end{lemma}
\begin{proof}
For any $x,y\in L^2(\mathbb R^d)$, consider their coordinate sequences
$\{x_n\},\{y_n\}\in \ell^2(\mathbb N_0^d)$ w.r.t.\ the basis $\{|n\rangle\}$. We have
\begin{gather*}
\vert\langle x,B_\alpha y\rangle\vert=
\bigg|\sum_{m,n}\overline{x_m} b^\alpha_{mn} y_n\bigg| \leq
\sum_{m,n}(1+|m-n|)^\alpha |x_m| |y_n|\\
\hphantom{\vert\langle x,B_\alpha y\rangle\vert=
\bigg|\sum_{m,n}\overline{x_m} b^\alpha_{mn} y_n\bigg|}{}
\leq\sum_{m,k}(1+|k|)^\alpha |x_m| |y_{k+m}|
\leq\sum_k(1+|k|)^\alpha \Vert x\Vert_{L^2} \Vert y\Vert_{L^2} ,
\end{gather*}
which concludes the proof.
\end{proof}

\begin{lemma}\label{anorm}
For any $\alpha>d$, there exists a constant $C_\alpha>0$ such that the following holds:
\[
\Vert\phi\Vert_a\leq C_\alpha \Vert\phi\Vert_{\infty,\alpha} ,\qquad \phi\in{\cal L}_{\infty,\alpha} .
\]
\end{lemma}

\begin{proof}
With notation as in the previous proof we have, for $\phi\in{\cal L}_{\infty,\alpha}$,
\begin{gather*}
\bigg|\sum_{m,n}|\phi_{mn}| \overline{x_n}\,y_m\bigg| \leq
\sum_{m,n}\big|b^\alpha_{mn} \phi_{mn}\big|
\big|b^{-\alpha}_{mn} \overline{x_n} y_m\big|
\leq \sup_{m,n}\big|b^\alpha_{mn}\phi_{mn}\big|
 \Vert B_{-\alpha}\Vert_{\rm op} \Vert x\Vert_{L^2} \Vert y\Vert_{L^2} ,
\end{gather*}
which evidently implies the claim by Lemma~\ref{Bbounded}.
\end{proof}

\begin{lemma}\label{trace}
For any $\alpha<-d$ there exists a constant $C'_\alpha>0$ such that
\[
\Vert\phi\Vert_{1,\alpha}\leq C'_{\alpha} \Vert\phi\Vert_1 ,\qquad \phi\in{\cal L}_1 ,
\]
where ${\cal L}_1$ denotes the space of trace class operators and
$\Vert\cdot\Vert_1$ is the standard trace-norm.
\end{lemma}

\begin{proof}
Writing the trace class operator $\phi$ as $\frac
12(\phi+\phi^*)+\frac i2(i\phi^*-i\phi)$ we may assume that $\phi$ is
self-adjoint. Then, writing $\phi= \phi_+-\phi_-$ where $\phi_\pm$ are positive
operators each of which has trace norm at most that of $\phi$, we can
assume $\phi$ is positive. In this case the Cauchy--Schwarz inequality
gives
\[
|\phi_{kl}|  \leq \big(\phi_{kk} \phi_{ll}\big)^{1/2},
\]
and hence
\begin{gather*}
\sum_{k,l}b_{kl}^{\alpha} |\phi_{kl}|\leq
\sum_{k,l}b_{kl}^{\alpha} \phi_{kk}^{1/2} \phi_{ll}^{1/2}
\leq \Vert B_{\alpha}\Vert_{\rm op} \sum_k\phi_{kk}
=\Vert B_{\alpha}\Vert_{\rm op} \Vert\phi\Vert_1 ,
\end{gather*}
which proves the claim thanks to Lemma~\ref{Bbounded}.
\end{proof}

\begin{lemma}\label{product}
$a)$  For any $\alpha\geq 0$ there exists a constant $C_{1,\alpha}>0$
such that
\begin{gather}\label{prodest1}
\Vert\phi \psi\Vert_{2,\alpha}\leq C_{1,\alpha}(\Vert\phi\Vert_{2,\alpha} \Vert\psi\Vert_a
+\Vert\phi\Vert_a \Vert\psi\Vert_{2,\alpha}),\qquad \phi, \psi\in {\cal L}_{2,\alpha} .
\end{gather}

$b)$  For any $\alpha> d$ there exists a constant $C_{2,\alpha}>0$
such that
\begin{gather}\label{prodest2}
\Vert\phi \psi\Vert_{1,\alpha}\leq C_{2,\alpha}(\Vert\phi\Vert_{2,2\alpha} \Vert\psi\Vert_2
+\Vert\phi\Vert_2 \Vert\psi\Vert_{2,2\alpha}),\qquad \phi,\psi\in {\cal L}_{2, 2\alpha} .
\end{gather}
\end{lemma}

\begin{proof}
First, note that for $\alpha\geq 0$
\begin{gather}
\label{Bsum}
b_{mn}^\alpha\leq c(\alpha)\,(b_{mk}^\alpha+b_{kn}^\alpha) ,
\qquad  m,n,k\in\mathbb N_0^d,
\end{gather}
where the constant $c(\alpha)$ depends only on $\alpha$.
We then have
\begin{gather*}
\Vert\phi \psi\Vert_{2,\alpha}^2 =\sum_{m,n}b^{2\alpha}_{mn}\bigg|\sum_k\phi_{mk}
\psi_{kn}\bigg|^2
 \leq
C\sum_{m,n}\bigg(\sum_kb^{\alpha}_{mk} |\phi_{mk}| \, |\psi_{kn}|
+\sum_kb^{\alpha}_{kn} |\phi_{mk}|\,|\psi_{kn}|\bigg)^2\\
\hphantom{\Vert\phi \psi\Vert_{2,\alpha}^2}{}
\leq
2C\sum_{m,n}\bigg(\sum_kb^{\alpha}_{mk} |\phi_{mk}|\,|\psi_{kn}|\bigg)^2
+2C\sum_{m,n}\bigg(\sum_kb^{\alpha}_{kn} |\phi_{mk}|\,|\psi_{kn}|\bigg)^2\\
\hphantom{\Vert\phi \psi\Vert_{2,\alpha}^2}{}
=2C\,\big(\Vert(U_{\alpha}\widetilde\phi) \widetilde\psi\Vert_2^2+
\Vert\widetilde\phi (U_{\alpha}\widetilde\psi)\Vert_2^2\big)\\
\hphantom{\Vert\phi \psi\Vert_{2,\alpha}^2}{}
\leq 2C \big(\Vert U_{\alpha}\widetilde\phi\Vert_2^2
\Vert\widetilde\psi\Vert_{\rm op}^2+
\Vert U_\alpha\widetilde\psi\Vert_2^2
\Vert\widetilde\phi\Vert_{\rm op}^2\big)
 \leq 2C \big(\Vert\phi\Vert_{2,\alpha}
\Vert\psi\Vert_a+\Vert\psi\Vert_{2,\alpha} \Vert\phi\Vert_a\big)^2 ,
\end{gather*}
where  $C= c(\alpha)^2$ and $U_\alpha:\mathcal H_\alpha\to\mathcal H_0$ is
the unitary operator def\/ined by (\ref{Ualpha}). This establishes~(\ref{prodest1}).

Using again~\eqref{Bsum}, we have
\begin{gather}
\Vert\phi \psi\Vert_{1,\alpha} =\sum_{m,n}b^\alpha_{mn}\bigg|\sum_k\phi_{mk}
\psi_{kn}\bigg| \leq
c(2\alpha)\sum_{m,n,k}b^{-\alpha}_{mn}\big(|b^{2\alpha}_{mk} \phi_{mk}|\,
|\psi_{kn}|+|\phi_{mk}|\,|b^{2\alpha}_{kn} \psi_{kn}|\big)\nonumber\\
\hphantom{\Vert\phi \psi\Vert_{1,\alpha}}{}
=c(2\alpha)\big(\Vert(U_{2\alpha}\widetilde\phi) \widetilde\psi\Vert_{1,-\alpha}+
\Vert\widetilde\phi (U_{2\alpha}\widetilde\psi)\Vert_{1,-\alpha}\big) .\label{34}
\end{gather}
 Lemma~\ref{trace} and a H\"older inequality now yield, for $\alpha>d$,
\[
\Vert(U_{2\alpha}\widetilde\phi) \widetilde\psi\Vert_{1,-\alpha}
\leq C'_{-\alpha}\Vert(U_{2\alpha}\widetilde\phi) \widetilde\psi\Vert_1
\leq C'_{-\alpha}\Vert U_{2\alpha}\widetilde\phi\Vert_2\Vert\psi\Vert_2
= C'_{-\alpha}\Vert\phi\Vert_{2,2\alpha}\Vert\psi\Vert_2 .
\]
Combining this with \eqref{34} gives \eqref{prodest2}.
\end{proof}

\section{Decay of free solutions}
\label{secth2}

The goal in this section is to prove Theorem \ref{th2}. The main
step in the proof is to establish appropriate time-decay estimates for the matrix elements
$\left(e^{-it{\bf\Delta}}\right)_{nm,kl}$ of the free time-evolution
operator with respect to the orthonormal basis
$\{(|n\rangle\langle m|)\}$ for ${\cal H}$. This is our f\/irst objective.

Since the Weyl map $W$ is unitary up to a constant factor by \eqref{Wunitary}, the matrix elements of the heat operator $e^{-t{\bf\Delta}}$, $t>0$, can be
obtained in closed form from the well-known  expression for
the heat kernel in Euclidean space and the relation \eqref{tria},
making use of the fact that the functions $W^{-1}(|n\rangle\langle
m|)$ can be computed explicitly in terms of Laguerre polynomials. The
result is given in \cite{Raimar} and reads, in case $d=1$,
\begin{gather*}
\left(e^{-t{\bf\Delta}}\right)_{nm,kl}=
\delta_{m+k,n+l}\sum_{v=0}^{\min\{m,l\}} C_{nm,kl,v}
 \frac{t^{m+l-2v}}{(1+t)^{m+k+d}} ,
\end{gather*}
where
\begin{gather}
\label{matconst}
C_{nm,kl,v}:=
\sqrt{{n\choose m-v}{k\choose l-v}
{m\choose m-v}{l\choose l-v}} ,
\end{gather}
for arbitrary non-negative integers $n$, $m$, $k$, $l$. By convention the
binomial coef\/f\/icient $n\choose m$ vanishes unless $0\leq m\leq n$. Note that the presence of
the Kronecker delta factor is a consequence of rotational invariance
of the heat kernel in two-dimensional space, which in the operator
formulation entails that $ e^{-t{\bf\Delta}}$ commutes with $e^{-isa^*a}$ for
$s\in \mathbb R$, where we have dropped the subscript on the
annihilation and creation operators when $d=1$.

Since ${\bf\Delta}$ is a positive,  self-adjoint operator we obtain
$\left(e^{-it{\bf\Delta}}\right)_{nm,kl}$ for $t\in\mathbb R$ by analytic
continuation in $t$, that is
\begin{gather}
\label{mateluni}
\left(e^{-it{\bf\Delta}}\right)_{nm,kl}=
\delta_{m+k,n+l}\sum_{v=0}^{\min\{m,l\}} C_{nm,kl,v}
 \frac{(it)^{m+l-2v}}{(1+it)^{m+k+d}} .
\end{gather}

We next observe that these  matrix
elements are expressible in terms of Jacobi polynomials.

\begin{lemma}
\label{jacobi}
For $d=1$ and $l\leq m,k\leq n$, we have
\begin{gather}\label{jacrep}
\left(e^{-it{\bf\Delta}}\right)_{nm,kl}=
\delta_{m+k,n+l}
\sqrt{\frac{n! l!}{m! k!}}
\frac{(it)^{m+l}}{(1+it)^{m+k+1}}\big(1+t^{-2}\big)^l
  P_{l}^{n-m,m-l}\left(\frac{t^2-1}{t^2+1}\right) ,
\end{gather}
where $P_l^{\alpha,\beta}(X)$, $l\in\mathbb N_0$, $\alpha,\beta\geq 0$,
$X\in[-1,1]$, are the classical Jacobi polynomials with standard
normalization~{\rm \cite{Szego}}.
\end{lemma}

\begin{proof}
For  $m+k=n+l$ and $n\geq m$ we have
\begin{gather*}
C_{nm,kl,v}
=\sqrt{\frac{n! l!}{m! k!}}{m\choose v}{l+n-m\choose l-v},
\end{gather*}
and consequently, for $l\leq m$,
\begin{gather}
\label{matelements}
\left(e^{-it{\bf\Delta}}\right)_{nm,kl}=
\delta_{m+k,n+l}\sqrt{\frac{n! l!}{m! k!}}
\frac{(it)^{m+l}}{(1+it)^{m+k+1}}
\sum_{v=0}^{l} {m\choose v}{l+n-m\choose l-v}
 \big(-t^{-2}\big)^{v} .
\end{gather}
Recall now \cite{Szego} that the classical Jacobi polynomials
$P_l^{\alpha,\beta}(X)$, $l\in\mathbb N_0,$ are orthogonal
w.r.t.\ the weight function
\begin{gather}\label{weight}
w(X)   =  (1-X)^\alpha(X+1)^\beta,\qquad X\in[-1,1],
\end{gather}
for f\/ixed values of $\alpha, \beta > -1$. For our purposes we may
restrict attention to integer values of $\alpha$ and  $\beta$. A  convenient
explicit form of $P_l^{\alpha,\beta}(X)$ with standard normalization is
\[
P_l^{\alpha,\beta}(X)=\sum_{j=0}^l{l+\alpha\choose l-j}{l+\beta\choose
  j}\left(\frac{X-1}{2}\right)^j\left(\frac{X+1}{2}\right)^{l-j} .
\]
With
\begin{gather}\label{X-t}
X=(t^2-1)/(t^2+1)\in [-1,1]
\end{gather}
this gives
\[
P_l^{\alpha,\beta}\left(\frac{t^2-1}{t^2+1}\right)
= \big(1+t^{-2}\big)^l\sum_{j=0}^l{l+\alpha\choose l-j}
{l+\beta\choose j}\big(-t^{-2}\big)^{j}
\]
and the result follows by comparison with~\eqref{matelements}.
\end{proof}

The representation~\eqref{jacrep} combined with rather recent
uniform estimates on the Jacobi polynomials may now be used to derive the
following bound on the matrix elements in question.

\begin{lemma}
\label{Uniest1}
For $d\geq 1$, there exists a constant $C_d$, independent of $n,m,k,l
\in\mathbb N_0^d$ and
$t\in\mathbb R$, such that
\begin{gather}\label{decay0}
\big|\left(e^{-it{\bf\Delta}}\right)_{nm,kl}\big| \leq  C_d |t|^{-\frac
d2} ,\qquad |t|\geq 1 .
\end{gather}
\end{lemma}

\begin{proof} From the def\/inition of ${\bf\Delta}$ it follows that the
  matrix elements in question factorize into those for the
  one-dimensional case, so it suf\/f\/ices to consider $d=1$.

 Since $e^{-it{\bf\Delta}}$ is unitary, it follows from
  \eqref{mateluni} that
 \[
\left(e^{-it{\bf\Delta}}\right)_{nm,kl} =  \left(e^{-it{\bf\Delta}}\right)_{kl,nm}.
\]
 Moreover, since
$W(\overline f)=W(f)^*$ and the heat kernel on $\mathbb R^2$ is
symmetric in its two arguments we likewise have
 \[
\left(e^{-it{\bf\Delta}}\right)_{nm,kl} =  \left(e^{-it{\bf\Delta}}\right)_{mn,lk} .
\]
These symmetries may also be checked directly from \eqref{mateluni} and \eqref{matconst}.

Taking into account that $m+k=n+l$ for non-vanishing matrix elements
we can therefore assume that $l\leq m,k\leq n$. Lemma~\ref{jacobi} then yields
\begin{gather}
\label{estimategene}
\left|\left(e^{-it{\bf\Delta}}\right)_{nm,kl}\right|
 = \delta_{m+k,n+l}
\sqrt{\frac{n! l!}{m! k!}}(1+t^2)^{-1/2}
\big(1+t^{-2}\big)^{(l-n)/2}|t|^{l-k}
\left|P_l^{\alpha,\beta}\left(\frac{t^2-1}{t^2+1}\right)\right|,
\end{gather}
where we have set
\begin{gather*}%\label{alphabeta}
\alpha:=n-m \qquad\mbox{and}\qquad \beta:=m-l .
\end{gather*}
Introducing the orthonormal Jacobi polynomials ${\bf
  P}_l^{\alpha,\beta}$ of degree $l\geq 0$ by normalizing w.r.t.\ the
$L^2$-norm def\/ined by the weight \eqref{weight}, one f\/inds (see e.g.~\cite[p.~67]{Szego})
\[
{\bf P}_l^{\alpha,\beta}(X)  =  \sqrt{\frac{(2l+\alpha +\beta +1)}{2^{\alpha
    +\beta +1}}\frac{(l+\alpha +\beta)! l!}{(l+\alpha)! (l+\beta)!}}
P_l^{\alpha,\beta}(X) ,
\]
and \eqref{estimategene} can be rewritten as
\begin{gather}
\left|\left(e^{-it{\bf\Delta}}\right)_{nm,kl}\right|
 = \delta_{m+k,n+l}\left(l+\frac{\alpha +\beta +1}2\right)^{-\frac 12}\nonumber\\
 \hphantom{\left|\left(e^{-it{\bf\Delta}}\right)_{nm,kl}\right|=}{}\times
\big(1+t^2\big)^{-1/2} (1-X)^{\alpha/2}(1+X)^{\beta/2}
\big|{\bf P}_l^{\alpha,\beta}(X)\big| ,\label{estimategene2}
\end{gather}
with $X$ given by \eqref{X-t}.

We now use the following uniform bound for the orthonormal Jacobi polynomials,
proven in~\cite{erdel}:
\begin{gather}\label{erdelyi}
(1-X)^{\alpha/2+1/4}
(1+X)^{\beta/2+1/4}\big|{\bf P}_l^{\alpha,\beta}(X)\big| \leq \sqrt{\frac{2 e}{\pi}}\sqrt{2+
\sqrt{\alpha^2+\beta^2}} ,
\end{gather}
valid for all $X\in]-1,1[$, $l\geq 0 $ and
$\alpha, \beta\geq-1/2$.

Applying this inequality in conjunction with \eqref{estimategene2} we
obtain, for $m,n,k,l\geq 0$ and   $|t|\geq 1$,
\begin{gather*}%\label{matest2}
\left|\left(e^{-it{\bf\Delta}}\right)_{mn,kl}\right|
 \leq C \delta_{m+k,n+l}
\left(\frac{2+
\sqrt{\alpha^2+\beta^2}}
{2l+\alpha +\beta +1}\right)^{\frac 12}
|t|^{-\frac 12}\leq C' \delta_{m+k,n+l} |t|^{-\frac 12} ,
\end{gather*}
which is the announced result.
\end{proof}
\setcounter{section}{4}
\setcounter{subsection}{0}
\setcounter{subsubsection}{0}
\setcounter{equation}{32}
% Original source lines 987--1031
We are now ready to prove Theorem~\ref{th2}.

\begin{proof}[Proof of Theorem~\ref{th2}]
Recalling the def\/inition \eqref{Ualpha} of $U_\alpha$ and that
$b_{mn}$ only depends on $|n-m|$, it follows from the presence of the
Kronecker-delta factor in $(e^{-it{\bf{\Delta}}})_{nm,kl}$ that
\[
e^{-it{\bf{\Delta}}_\alpha}(|n\rangle\langle m|) = U_\alpha^{-1}
e^{-it{\bf{\Delta}}}U_\alpha(|n\rangle\langle m|) = e^{-it{\bf{\Delta}}}(|n\rangle\langle m|)
\]
is independent of $\alpha$. Since
$\{b_{mn}^{-\alpha}|n\rangle\langle m|\}$ is an orthonormal basis
for ${\cal H}_\alpha$ we obtain, for $\alpha\geq 0$ and
\[
\phi =
\sum_{m,n}\phi_{mn}|n\rangle\langle m| \in {\cal H}_\alpha\subseteq
{\cal H} ,
\]
that
\[
e^{-it{\bf{\Delta}}_\alpha}\phi =
\sum_{n,m}\phi_{mn}e^{-it{\bf{\Delta}}_\alpha}(|n\rangle\langle m|) =
\sum_{n,m}\phi_{mn}e^{-it{\bf{\Delta}}}(|n\rangle\langle m|) = e^{-it{\bf{\Delta}}}\phi ,
\]
i.e. $e^{-it{\bf{\Delta}}_\alpha}$ equals the restriction of $e^{-it{\bf{\Delta}}}$
to ${\cal H}_\alpha$.
Applying  Lemma~\ref{anorm} we then get, for $\beta>d$ and
$\phi\in {\cal H}_\alpha\cap {\cal L}_{1,\beta}$,
\begin{gather*}
\big\Vert e^{-it{\bf{\Delta}}_\alpha}\phi\big\Vert_a
 \leq
 C_\beta \sup_{n,m}\bigg|b_{nm}^{\beta}
\sum_{k,l}\big(e^{-it{\bf{\Delta}}}\big)_{nm,kl}\,\phi_{kl}\bigg|
 \leq C_\beta \sup_{n,m}\sum_{k,l}\Big|
\big(e^{-it{\bf{\Delta}}}\big)_{nm,kl} b_{kl}^{\beta} \phi_{kl}\Big|\\
\phantom{\big\Vert e^{-it{\bf{\Delta}}_\alpha}\phi\big\Vert_a}{}
\leq
C_\beta \sup_{n,m,k,l}\Big|\big(e^{-it{\bf{\Delta}}}\big)_{nm,kl}\Big| \sum_{k,l}\big|b_{kl}^{\beta} \phi_{kl}\big|
= C_\beta \sup_{n,m,k,l}\Big|\big(e^{-it{\bf{\Delta}}}\big)_{nm,kl}\Big|
\Vert\phi\Vert_{1,\beta} ,
\end{gather*}
where, in the second step, we have once more made use of the
Kronecker-delta factor in $(e^{-it{\bf{\Delta}}})_{nm,kl}$.
The claim now follows from Lemma~\ref{Uniest1}.
\end{proof}
\setcounter{section}{4}
\setcounter{subsection}{0}
\setcounter{subsubsection}{0}
\setcounter{equation}{32}
% Original source lines 1062--1186
\section{Existence of wave operators for small data}
\label{secth3}

\begin{proof}[Proof of Theorem~\ref{th3}.]
~~Both statements follow from~ \cite[Theorem 16]{Reed} once we verify
the following four conditions for $\alpha> 2d$ and some $\delta>0$,
where $p$ denotes the lowest degree of monomials occurring in $F$:
\begin{itemize}\itemsep=0pt
\item[i)] There exists a constant $c_1>0$ such that
\[
\Vert e^{-it{\bf{\Delta}}_\alpha}\phi\Vert_a\leq
c_1 |t|^{-d/2} \Vert\phi\Vert_{1,\alpha/2},\qquad |t|\geq 1 .
\]
\item[ii)] There exists a constant $c_2>0$ such that
\[
\Vert\phi\Vert_a\leq c_2 \Vert\phi\Vert_{2,\alpha} .
\]
\item[iii)] There  exists a constant $c_3>0$ such that
\[
\Vert\phi F(\phi^*\phi)-\psi F(\psi^*\psi)\Vert_{2,\alpha}
\leq c_3\big(\Vert\phi\Vert_a+\Vert\psi\Vert_a\big)^{2p-1}
\Vert\phi-\psi\Vert_{2,\alpha} ,\qquad \mbox{if}\ \ \Vert\phi\Vert_{2,\alpha},
\Vert\psi\Vert_{2,\alpha}\leq\delta .
\]
\item[iv)]  There  exists a constant $c_4>0$ such that
\begin{gather*}
\Vert\phi F(\phi^*\phi)-\psi F(\psi^*\psi)\Vert_{1,\alpha/2}\\
\qquad {} \leq c_4\Big\{\big(\Vert\phi\Vert_a+\Vert\psi\Vert_a\big)^{2p-2}
\Vert\phi-\psi\Vert_a
  +
\big(\Vert\phi\Vert_a+\Vert\psi\Vert_a\big)^{2p-1}
\Vert\phi-\psi\Vert_{2,\alpha}\Big\} ,\\
\qquad \qquad \mbox{if} \ \Vert\phi\Vert_{2,\alpha},
\Vert\psi\Vert_{2,\alpha}\leq\delta .
\end{gather*}
\end{itemize}
Moreover, it is required that the constants $c_3$, $c_4$ can be chosen
arbitrarily small by choosing $\delta$ small enough.

To be specif\/ic, the stated basic assumptions about $F$ ensure that
 $2p-1\geq 1$ and $\frac d2(2p-1)> 1$ for all
$d\geq 1$, which implies \eqref{solsmall1} by  \cite[Theorem
16]{Reed}. If, in addition, $F$ has no linear term, we have $2p-1>1$
and \eqref{solsmall2} follows similarly.

Condition i) above is a particular case of Theorem~\ref{th2}. That
condition ii) holds, follows from
\begin{gather*}%\label{alessalpha}
\Vert\phi\Vert_a^2 \leq \Vert\widetilde\phi\Vert^2_2 = \Vert\phi\Vert_{2,0}^2 \leq
\sum_{m,n}b_{mn}^{2\alpha}\big|\phi_{mn}\big|^2 = \Vert\phi\Vert^2_{2,\alpha},\qquad \alpha\geq0,
\end{gather*}
where we have used that the Hilbert--Schmidt norm dominates the operator
norm in the f\/irst step and that $b^{2\alpha}_{mn}\geq 1$ for
$\alpha\geq 0$ in the third step.

To establish iii) and iv) we make use of Lemma~\ref{product}. As a
consequence of the inequality
$\Vert\phi\,\psi\Vert_a\leq\Vert\phi\Vert_a \Vert\psi\Vert_a$, which
follows immediately from
\[
(\widetilde{\phi\psi})_{mn} = \bigg|\sum_k\phi_{mk} \psi_{kn}\bigg| \leq \sum_k|\phi_{mk}|\,|\psi_{kn}| = (\widetilde\phi \widetilde\psi)_{mn} ,
\]
the f\/irst part of the Lemma implies
\begin{gather}
\label{prodnorm1}
\Vert\phi_1\cdots\phi_r\Vert_{2,\alpha}\leq
C(\alpha,r)  \sum_{i=1}^r\Vert\phi_i\Vert_{2,\alpha} \prod_{j\ne i}\Vert\phi_j\Vert_a ,
\end{gather}
valid for $\alpha\geq 0$ and some constant $C(\alpha,r)$ depending
on $\alpha$ and $r$ only.

The second statement of Lemma~\ref{product} implies, for $\alpha>2d$,
\begin{gather*}%\label{prodnorm2a}
\Vert\phi_1\cdots\phi_r\Vert_{1,\alpha/2} \leq
C_{2,\alpha/2} \bigg(\Vert\phi_1\Vert_{2,\alpha} \bigg\Vert\prod_{i=2}^r\phi_i\bigg\Vert_2 +
 \Vert\phi_1\Vert_2 \bigg\Vert\prod_{i=2}^r\phi_i\bigg\Vert_{2,\alpha}\bigg)\\
 \hphantom{\Vert\phi_1\cdots\phi_r\Vert_{1,\alpha/2}}{}
 \leq
2C_{2,\alpha/2} \Vert\phi_1\Vert_{2,\alpha} \bigg\Vert\prod_{i=2}^r\phi_i\bigg\Vert_{2,\alpha}
\end{gather*}
 where, in the second step, it has been used that $\Vert\phi\Vert_{2,\alpha}$ is an
increasing function of $\alpha$.

The last expression can now be estimated by making use of~\eqref{prodnorm1}. Thus we obtain
\begin{gather}\label{prodnorm2}
\Vert\phi_1\cdots\phi_r\Vert_{1,\alpha/2}
\leq 2 C(\alpha,r-1)C_{2,\alpha/2} \Vert\phi_1\Vert_{2,\alpha} \sum_{i=2}^r\Vert\phi_i\Vert_{2,\alpha}
\prod_{j\ne 1,i}\Vert\phi_j\Vert_a .
\end{gather}
Now, let $p$ be an arbitrary positive integer and write
\begin{gather*}
\phi |\phi|^{2p}-\psi |\psi|^{2p}
=\sum_{i=0}^{p-1}\phi |\phi|^{2(p-1-i)}
(\phi^* (\phi-\psi)+(\phi^*-\psi^*) \psi) |\psi|^{2i}
+(\phi-\psi) |\psi|^{2p} .
\end{gather*}
Since the norms $\Vert\cdot\Vert_{2,\alpha}$, $\Vert\cdot\Vert_a$, $\Vert\cdot\Vert_{1,\alpha}$
are  $*$-invariant, an application of~\eqref{prodnorm1}
and~\eqref{prodnorm2} yields the inequalities
\begin{gather*}
\Vert\phi |\phi|^{2p}-\psi |\psi|^{2p}\Vert_{2,\alpha}
\leq c'_3 \big(\Vert\phi\Vert_{2,\alpha}+\Vert\psi\Vert_{2,\alpha}\big)
\big(\Vert\phi\Vert_a+\Vert\psi\Vert_a\big)^{2p-1}\Vert\phi-\psi\Vert_{2,\alpha} ,
\end{gather*}
and
\begin{gather*}
 \Vert\phi |\phi|^{2p}-\psi |\psi|^{2p}\Vert_{1,\alpha/2}
 \leq c'_4
\Big\{\big(\Vert\phi\Vert_{2,\alpha}+\Vert\psi\Vert_{2,\alpha}\big)^2
\big(\Vert\phi\Vert_a+\Vert\psi\Vert_a\big)^{2p-2}\Vert\phi-\psi\Vert_a\\
\phantom{\Vert\phi |\phi|^{2p}-\psi |\psi|^{2p}\Vert_{1,\alpha/2}\leq}{}
+
\big(\Vert\phi\Vert_{2,\alpha}+\Vert\psi\Vert_{2,\alpha}\big)
\big(\Vert\phi\Vert_a+\Vert\psi\Vert_a\big)^{2p-1}\Vert\phi-\psi\Vert_{2,\alpha}\Big\} ,
\end{gather*}
respectively, for suitable constants $c'_3$, $c'_4$, where the
obvious inequality
\[
\Vert\phi\Vert_a \leq \Vert\phi\Vert_{2,\alpha} ,\qquad \alpha\geq 0 ,
\]
has also been used. This evidently proves iii) and iv), with $c_3$ and
$c_4$ of order $\delta$, if $F$ is a~monomial of degree $p$. The same bounds then follow immediately for an
arbitrary polynomial~$F$ without constant term, if~$p$ denotes the
lowest degree of monomials occurring in~$F$.
\end{proof}
\begin{thebibliography}{99}
\bibitem[2]{Caz}
Cazenave T.,
Semilinear Schr{\"o}dinger equations,
{\it Courant Lecture Notes in Mathematics}, Vol.~10, Amer. Math. Soc., Providence, RI, 2003.

\bibitem[4]{DJN}
Durhuus B., Jonsson T., Nest R.,
The existence and stability of noncommutative scalar solitons,
\href{http://dx.doi.org/10.1007/s00220-002-0721-4}{{\it Comm. Math. Phys.}} {\bf 233} (2003), 49--78,
\href{http://arxiv.org/abs/hep-th/0107121}{hep-th/0107121}.

\bibitem[5]{erdel}
Erd{\'e}lyi T., Magnus  A.P., Nevai P.,
Generalized Jacobi weights, Christof\/fel functions, and Jacobi polynomials,
\href{http://dx.doi.org/10.1137/S0036141092236863}{{\it SIAM J. Math. Anal.}} {\bf 25} (1994), 602--614.

\bibitem[6]{GGBISV}
Gayral V., Gracia-Bond\'{\i}a J.M., Iochum B., Sch\"ucker T.,  V\'arilly J.C.,
Moyal planes are spectral triples,
\href{http://dx.doi.org/10.1007/s00220-004-1057-z}{{\it Comm. Math. Phys.}} \textbf{246} (2004), 569--623,
\href{http://arxiv.org/abs/hep-th/0307241}{hep-th/0307241}.

\bibitem[7]{Ginibre}
Ginibre J.,
An introduction to nonlinear Schr{\"o}dinger equations,
in Nonlinear Waves (Sappore 1995), Editors R.~Agemi, Y.~Giga and T.~Ozawa,
{\it GAKUTO Internat. Ser. Math. Sci. Appl.}, Gakk{\=o}tosho, Tokyo, 1997,  85--133.

\bibitem[8]{GiVe}
Ginibre J., Velo G.,
Time decay of f\/inite energy solutions of the nonlinear Klein--Gordon and Schr{\"o}dinger equations,
{\it Ann. Inst. H. Poincar{\'e} Phys. Th{\'e}or.} {\bf 43} (1985), 399--442.

\bibitem[9]{GMS}
Gopakumar R., Minwalla S., Strominger A.,
Noncommutative solitons,
\href{http://dx.doi.org/10.1088/1126-6708/2000/05/020}{{\it J. High Energy Phys.}} {\bf 2000}, (2000), no.~5, 020, 27~pages,
\href{http://arxiv.org/abs/hep-th/0003160}{hep-th/0003160}.

\bibitem[11]{HKL}
Harvey J., Kraus P., Larsen F.,
Exact noncommutative solitons,
 \href{http://dx.doi.org/10.1088/1126-6708/2000/12/024}{{\it J. High Energy Phys.}} {\bf 2000}, (2000), no.~12, 024, 24~pages,
 \href{http://arxiv.org/abs/hep-th/0010060}{hep-th/0010060}.

\bibitem[13]{Reed}
Reed M.,
Abstract non-linear wave equations,
{\it Lecture Notes in Mathematics}, Vol.~507, Springer-Verlag,
  Berlin~-- New York, 1976.

\bibitem[14]{RS}
Reed  M., Simon B.,
Methods of modern mathematical physics, Vol.~II, Academic Press, New York~-- London, 1975.

\bibitem[16]{Szego}
Szeg\"o G.,
Orthogonal polynomials,  {\it American Mathematical Society Colloquium Publications}, Vol.~23, Amer. Math. Soc., Providence, RI, 1959.

\bibitem[17]{Strauss} 
Strauss W.A.,
Nonlinear wave equations,
{\it CBMS Regional Conference Series in Mathematics}, Vol.~73, Amer. Math. Soc., Providence, RI, 1989.

\bibitem[18]{Raimar} Wulkenhaar R., Renormalisation of noncommutative $\varphi_4^4$-theory to all orders, Habilitation Thesis, University of Vienna, 2003, available at \url{http://www.math.uni-muenster.de/u/raimar/}.

\end{thebibliography}
\end{document}
